\begin{afterword}
\summaryhead

The post opens with an argument for keeping children's belief in Santa Claus: it gives them wonder and makes them behave. The author calls this a false dilemma. Even if believing in Santa beats doing nothing, it need not be the best option. Space launches and science fiction can supply wonder, and praise without bribes, the post says, produces good behavior whether or not adults are watching. Noble Lies, it concludes, are generally package deals, and a Third Alternative can supply the benefit without the lie.

Such false dilemmas persist, the post argues, because a search for a good policy stops as soon as it finds one that is defensible and convenient, and continues when the options found are costly to the searcher. The reader is told to be wary of calling a policy defensible rather than optimal, to ask whether a better option would be welcome, and to spend five minutes by the clock looking for one.

\responsehead

The post's logical point is correct. Showing that a policy beats doing nothing does not show that it beats every alternative, and the post names the error with its standard names. Its psychological claim, that a search ends as soon as a convenient option appears, is offered without evidence, but research supports something close, and a later post by the author relates it to Gilovich's ``motivated skepticism.'' The weaknesses lie in how far the post generalizes, and in its advice.

The generalizations are based on one example. The only Noble Lie examined is Santa Claus, an argument in quotation marks with no speaker. From it the post concludes that Noble Lies are ``generally'' package deals, that ``we can construct'' an alternative when the gain is needed, and that Noble Liars keep the lie they started with. No Noble Liar is named. ``Most people fail'' at deciding to look and deciding to accept comes with no evidence either. Even the example's own factual premise is stated too strongly: research on rewards supports the direction of the claim about bribes and praise, but not ``unconditional good behavior.''

The advice conflicts with the post's own premises. The post grants, with its house example, that a search among many options needs a rule for stopping short of the best. It then warns against defending a policy as ``defensible'' rather than ``optimal,'' and quotes ``The best is the enemy of the good'' in a new sense, although the proverb usually advises against holding out for the best. The warning needs a way to tell a motivated stop from a reasonable one. The first test offered is introspective: would you feel ``a tiny flash of reluctance''? Three paragraphs earlier the post said that the motives in question are ``subconscious influences,'' and does not say why asking would reveal them. The second test, five minutes by a physical clock, does not depend on introspection and can be followed as written.

In short: a correct point about false dilemmas, generalized from one unattributed example, with a first self-test that relies on introspection about motives the post calls subconscious.
\end{afterword}
